\section*{الفصل \ref{ch:g12:stereochemistry} --- الكيمياء الفراغية: الكيرالية والمتماكبات}

\begin{solution}{exo:g12:stereochemistry:1}
بوتان-2-ول: الكربون 2 (\ce{H}، \ce{OH}، \ce{CH3}، \ce{CH2CH3}).
بروبان-2-ول: لا شيء (فالكربون 2 يحمل مجموعتي \ce{CH3} متطابقتين).
2-كلورو بوتان: الكربون 2. حمض اللبنيك: الكربون 2 (\ce{H}، \ce{OH}،
\ce{CH3}، \ce{COOH}).
\end{solution}

\begin{solution}{exo:g12:stereochemistry:2}
الأول: ذرتا \ce{Cl} (لهما الأولوية على \ce{H}) في الجهة نفسها: Z.
الثاني: على اليسار، الأولوية للمجموعة \ce{CH3}، في الجهة العليا؛ وعلى اليمين،
الأولوية للمجموعة \ce{CH2CH3}، في الجهة السفلى: E.
\end{solution}

\begin{solution}{exo:g12:stereochemistry:3}
\omterm{def:g12:stereochemistry:chiral}{كيرالي}: القفاز والبرغي. غير \omterm{def:g12:stereochemistry:chiral}{كيرالي}: الملعقة والمكعب (فلكل منهما
مستوى مرآة). و\ce{CHFClBr} \omterm{def:g12:stereochemistry:chiral}{كيرالي} (أربع \omterm{def:g7:atoms-and-molecules:atom}{ذرات} مختلفة على الكربون)؛
أما \ce{CH2Cl2} فلا.
\end{solution}

\begin{solution}{exo:g12:stereochemistry:4}
\omterm{def:g6:pure-substances-and-mixtures:mixture}{خليط} من كميتين متساويتين من متماكبين مرآويين. والكارفون الراسيمي يحوي
المتماكبين المرآويين معًا، فيتلقى الأنف الرائحتين في آن واحد.
\end{solution}

\begin{solution}{exo:g12:stereochemistry:5}
لا: فالكربون 1 يحمل ذرتين متطابقتين، \ce{H} و\ce{H}؛ وتبديلهما
يعطي \omterm{def:g7:atoms-and-molecules:molecule}{الجزيء} نفسه.
\end{solution}

\begin{solution}{exo:g12:stereochemistry:6}
\setchemfig{atom sep=2em}
\chemfig{C(-[2]CH_2CH_3)(-[4]H_3C)(<[:-30]OH)(<:[:-150]H)} \quad مرآة
\quad \chemfig{C(-[2]CH_2CH_3)(-[0]CH_3)(<[:-150]HO)(<:[:-30]H)}
\end{solution}

\begin{solution}{exo:g12:stereochemistry:7}
2-كلورو بوتان: 2 (كربون واحد غير متناظر). 2,3-ثنائي كلورو بوتان: 3 (ذرتا
كربون غير متناظرتين ونصفان متماثلان: زوج من \omterm{def:g12:stereochemistry:enantiomer}{المتماكبات المرآوية} وشكل
ميزو). بنت-2-ين: 2 (Z وE).
\end{solution}

\begin{solution}{exo:g12:stereochemistry:8}
\omterm{def:g12:stereochemistry:enantiomer}{متماكبان مرآويان}؛ متماكبان لامرآويان؛ متماكبان لامرآويان؛ \omterm{def:g7:atoms-and-molecules:molecule}{الجزيء} نفسه (فحمض الطرطريك
ميزو ينطبق على صورته المرآوية).
\end{solution}

\begin{solution}{exo:g12:stereochemistry:9}
في \omterm{def:g12:stereochemistry:conformation}{التشكّل} المتعاقب الذي تتقابل فيه مجموعتا \ce{CH3}، تكون المجموعات
الكبيرة أبعد ما يمكن بعضها عن بعض: فهو الأكثر استقرارًا. وفي
\omterm{def:g12:stereochemistry:conformation}{التشكّل} المكسوف، تتواجه مجموعتا \ce{CH3} ويزاحم بعضهما بعضًا.
\end{solution}

\begin{solution}{exo:g12:stereochemistry:10}
نصفاه، وكل منهما \ce{CH(OH)COOH}، كلٌّ منهما صورة الآخر في المرآة: إذ يقسم
مستوى بين ذرتي الكربون المركزيتين \omterm{def:g7:atoms-and-molecules:molecule}{الجزيء} إلى نصفين كلٌّ منهما صورة
الآخر في المرآة. \omterm{def:g7:atoms-and-molecules:molecule}{والجزيء} الذي له مستوى مرآة هو صورة نفسه في المرآة: فهو غير
\omterm{def:g12:stereochemistry:chiral}{كيرالي}.
\end{solution}

\begin{solution}{exo:g12:stereochemistry:11}
\setchemfig{atom sep=1.7em}
(Z)-بنت-2-ين \chemfig{C(-[:150]H_3C)(-[:210]H)=C(-[:30]CH_2-CH_3)(-[:-30]H)}
و(E)-بنت-2-ين \chemfig{C(-[:150]H_3C)(-[:210]H)=C(-[:30]H)(-[:-30]CH_2-CH_3)}.
\end{solution}

\begin{solution}{exo:g12:stereochemistry:12}
كما في حمض الطرطريك، مع \ce{Br} بدل \ce{OH} و\ce{CH3}
بدل \ce{COOH}: ذرتا \ce{Br} كلتاهما بإسفين ممتلئ، أو كلتاهما بإسفين متقطع (زوج من
\omterm{def:g12:stereochemistry:enantiomer}{المتماكبات المرآوية})، وواحدة بإسفين ممتلئ والأخرى بمتقطع، ولها مستوى مرآة: وهي
شكل الميزو. ثلاثة \omterm{def:g12:stereochemistry:stereoisomer}{متماكبات فراغية}.
\end{solution}

\begin{solution}{exo:g12:stereochemistry:13}
النصف. يستطيع الكيميائي أن يفصل المتماكبين المرآويين (كما فعل باستور، لكن
بوسائل أخرى)، أو ألا يصنع إلا المتماكب المرآوي المفيد، بتفاعل
هو نفسه \omterm{def:g12:stereochemistry:chiral}{كيرالي} (حفّاز \omterm{def:g12:stereochemistry:chiral}{كيرالي}، أو إنزيم، أو \omterm{def:g5:raw-materials-and-recycling:raw-material}{مادة أولية}
كيرالية).
\end{solution}

\begin{solution}{exo:g12:stereochemistry:14}
ثلاثة: (2E,4E) و(2Z,4Z) و(2E,4Z)، والأخير هو \omterm{def:g7:atoms-and-molecules:molecule}{الجزيء} نفسه الذي هو
(2Z,4E) مقروءًا من الطرف الآخر.
\end{solution}

\begin{solution}{exo:g12:stereochemistry:15}
$2^3 = 8$ \omterm{def:g12:stereochemistry:stereoisomer}{متماكبات فراغية}، أي 4 أزواج من \omterm{def:g12:stereochemistry:enantiomer}{المتماكبات المرآوية}. ولكل
متماكب فراغي متماكب مرآوي واحد و$8 - 2 = 6$ \omterm{def:g12:stereochemistry:diastereomer}{متماكبات لامرآوية}.
\end{solution}

\begin{solution}{pb:g12:stereochemistry:1}
\textbf{1.} \ce{C4H6O6}؛ $M = 4 \times 12.0 + 6 \times 1.0 + 6 \times 16.0
= \qty{150.0}{g/mol}$.

\textbf{2.} مجموعتا كربوكسيل (حمض) ومجموعتا هيدروكسيل
(\omterm{def:g11:functional-groups:alcohol}{كحول}).

\textbf{3.} ذرتا الكربون المركزيتان، وكل منهما مرتبطة \omterm{def:g7:atoms-and-molecules:atom}{بالذرة} \ce{H} وبالمجموعات \ce{OH}
و\ce{COOH} و\ce{CH(OH)COOH}.

\textbf{4.} $2^2 = 4$.

\textbf{5.} كلتاهما بإسفين ممتلئ.

\textbf{6.} مجموعتا \ce{OH} بإسفينين متقطعين؛ ولا يمكن تطبيقه على
L: فهو حمض الطرطريك D، متماكبه المرآوي.

\textbf{7.} إذا أُدير حول الرابطة المركزية بحيث تتجه مجموعتا \ce{OH}
الوجهة نفسها، صار \omterm{def:g7:atoms-and-molecules:molecule}{للجزيء} مستوى بين ذرتي الكربون المركزيتين
يعكس أحد النصفين على الآخر.

\textbf{8.} لا: إنه حمض الطرطريك ميزو.

\textbf{9.} من التركيبات الأربع (ممتلئ، ممتلئ)، (متقطع، متقطع)،
(ممتلئ، متقطع)، (متقطع، ممتلئ)، الأخيرتان هما \omterm{def:g7:atoms-and-molecules:molecule}{الجزيء} نفسه،
شكل الميزو، لأن نصفي \omterm{def:g7:atoms-and-molecules:molecule}{الجزيء} متماثلان.

\textbf{10.} L وD: \omterm{def:g12:stereochemistry:enantiomer}{متماكبان مرآويان}. L وميزو: متماكبان لامرآويان.

\textbf{11.} عند \qty{169}{\celsius} أيضًا: فللمتماكبات المرآوية الخواص
الفيزيائية نفسها.

\textbf{12.} لا: فهو متماكب لامرآوي لحمض الطرطريك L، أي مركب مختلف
له خواصه الخاصة.

\textbf{13.} مستقبلات الأنف كيرالية وتميّز بين
الصورتين المرآويتين؛ أما مقياس الحرارة، غير \omterm{def:g12:stereochemistry:chiral}{الكيرالي}، فلا يستطيع.

\textbf{14.} لا، فهو \omterm{def:g6:pure-substances-and-mixtures:mixture}{خليط} من مركبين؛ وهو في مجمله غير
نشيط ضوئيًا، إذ يلغي كل من المتماكبين المرآويين أثر الآخر.

\textbf{15.} متماكب مرآوي واحد فقط: فكل بلورة مبنية من L أو من D
وحده.

\textbf{16.} \qty{5.0}{g} من كل منهما.

\textbf{17.} شكل الميزو مركب مختلف، ليس جزءًا من
\omterm{def:g12:stereochemistry:enantiomer}{الخليط الراسيمي} المكوّن من L وD.

\textbf{18.} بإذابة كل كومة: فالمحلولان يديران الضوء المستقطب
بالمقدار نفسه في اتجاهين متعاكسين (والبلورات كلٌّ منها صورة الأخرى
في المرآة).

\textbf{19.} أربعة: فذرتا الكربون غير المتناظرتين تحملان مجموعات مختلفة
(\ce{Cl} على إحداهما، و\ce{OH} على الأخرى)، فلا تركيبة لها مستوى
مرآة، ولا تكون تركيبتان الجزيءَ نفسه.

\textbf{20.} ثلاثة: L وD وميزو.
\end{solution}
